\documentclass[11pt,a4paper]{article}

% ---- Packages ----
\usepackage[utf8]{inputenc}
\usepackage[T1]{fontenc}
\usepackage{amsmath,amssymb,amsthm}
\usepackage{mathtools}
\usepackage[margin=1in,headheight=14pt]{geometry}
\usepackage{hyperref}
\usepackage{enumitem}
\usepackage{xcolor}
\usepackage{tcolorbox}
\usepackage{fancyhdr}
\usepackage{booktabs}

% ---- Hyperref ----
\hypersetup{
    colorlinks=true,
    linkcolor=red,
    citecolor=blue,
    urlcolor=blue
}

% ---- Header ----
\pagestyle{fancy}
\fancyhf{}
\fancyhead[L]{\textit{Lesson 6: Metric and Normed Spaces}}
\fancyhead[R]{\thepage}
\renewcommand{\headrulewidth}{0.4pt}

% ---- Theorem Environments ----
\theoremstyle{definition}
\newtheorem{definition}{Definition}[section]
\newtheorem{example}{Example}[section]
\newtheorem{exercise}{Exercise}[section]

\theoremstyle{plain}
\newtheorem{theorem}{Theorem}[section]
\newtheorem{lemma}[theorem]{Lemma}
\newtheorem{proposition}[theorem]{Proposition}
\newtheorem{corollary}[theorem]{Corollary}

\theoremstyle{remark}
\newtheorem{remark}{Remark}[section]

% ---- Colored Boxes ----
\tcbuselibrary{theorems,skins,breakable}

\newtcolorbox{keyresult}[1][]{
    colback=green!5!white,
    colframe=green!50!black,
    fonttitle=\bfseries,
    title=#1,
    breakable
}

\newtcolorbox{rlconnection}[1][]{
    colback=blue!5!white,
    colframe=blue!50!black,
    fonttitle=\bfseries,
    title=#1,
    breakable
}

\newtcolorbox{warning}[1][]{
    colback=red!5!white,
    colframe=red!50!black,
    fonttitle=\bfseries,
    title=#1,
    breakable
}

% ---- Operators ----
\newcommand{\R}{\mathbb{R}}
\newcommand{\N}{\mathbb{N}}
\newcommand{\Q}{\mathbb{Q}}
\newcommand{\Z}{\mathbb{Z}}
\newcommand{\C}{\mathbb{C}}
\newcommand{\Prob}{\mathbb{P}}
\newcommand{\E}{\mathbb{E}}
\newcommand{\indicator}{\mathbf{1}}
\DeclareMathOperator{\diam}{diam}
\DeclareMathOperator{\Int}{int}
\DeclareMathOperator{\cl}{cl}
\DeclareMathOperator{\Span}{span}
\DeclareMathOperator{\id}{id}

% ---- Title ----
\title{\textbf{Lesson 6: Metric and Normed Spaces}\\[10pt]
\large Metric Spaces, Banach Spaces, and Bounded Linear Operators\\[5pt]
\normalsize Session 6 of 102 $\mid$ Phase I: Mathematical Foundations $\mid$ 8\,h\\
\normalsize Primary text: Kreyszig, \emph{Introductory Functional Analysis with Applications}, Chapters 1--2}
\author{Curriculum: A Rigorous Learning Path to Multi-Agent Deep RL}
\date{}

\begin{document}

\maketitle

\begin{abstract}
This lesson develops the theory of metric spaces, normed spaces, and Banach spaces
from first principles. These are the natural settings for the function spaces and
operators that appear throughout reinforcement learning: value functions live in
Banach spaces, the Bellman operator acts on these spaces, and the convergence of
iterative algorithms (value iteration, policy iteration, TD-learning) relies on
completeness and contraction properties in these spaces. We prove every result,
including the equivalence of norms in finite dimensions, the characterization of
completeness via absolutely convergent series, the Arzel\`{a}--Ascoli theorem,
Riesz's lemma, and the Baire category theorem. The reader who masters this
material will have the tools to understand why the Bellman operator contracts,
why value iteration converges, and why infinite-dimensional spaces behave
fundamentally differently from $\R^n$.
\end{abstract}

\tableofcontents
\newpage

%% ============================================================
%% SECTION 1: INTRODUCTION
%% ============================================================
\section{Why Functional Analysis?}
\label{sec:intro}

Reinforcement learning algorithms operate on \emph{function spaces}. A value
function $V \colon \mathcal{S} \to \R$ assigns a real number to every state; a
Q-function $Q \colon \mathcal{S} \times \mathcal{A} \to \R$ assigns a value to
every state-action pair. When we write the Bellman optimality equation
\begin{equation}
\label{eq:bellman_intro}
V^*(s) = \max_{a \in \mathcal{A}} \left[ r(s,a) + \gamma \sum_{s'} p(s' \mid s,a)\, V^*(s') \right],
\end{equation}
we are asserting that $V^*$ is a \emph{fixed point} of an operator $T$ acting on
a space of functions. The fundamental questions are:

\begin{enumerate}[nosep]
    \item Does a fixed point exist?
    \item Is it unique?
    \item Does the iteration $V_{n+1} = T V_n$ converge to $V^*$?
    \item How fast does it converge?
\end{enumerate}

Answering these questions requires the language of functional analysis:

\begin{itemize}[nosep]
    \item \textbf{Metric spaces} (Section~\ref{sec:metric}) provide the notion of
    distance needed to say ``$V_n$ converges to $V^*$.''
    \item \textbf{Normed spaces} (Section~\ref{sec:normed}) give us a notion of
    ``size'' for functions and a way to measure error $\|V_n - V^*\|$.
    \item \textbf{Completeness} (Section~\ref{sec:banach}) ensures that Cauchy
    sequences converge, which is essential for the Banach fixed-point theorem.
    \item \textbf{Bounded linear operators} (Section~\ref{sec:operators}) formalize
    the maps between function spaces that arise in approximation theory.
    \item \textbf{Compactness} (Section~\ref{sec:compactness}) is the key to
    existence results and understanding the behavior of infinite-dimensional spaces.
    \item \textbf{The Baire category theorem} (Section~\ref{sec:baire}) underpins
    the three pillars of linear functional analysis: the Uniform Boundedness
    Principle, the Open Mapping Theorem, and the Closed Graph Theorem.
\end{itemize}

\begin{rlconnection}[The Bellman Operator as Motivation]
The Bellman optimality operator $T$ defined by the right-hand side of
\eqref{eq:bellman_intro} acts on the Banach space $(b(\mathcal{S}), \|\cdot\|_\infty)$
of bounded real-valued functions on the state space, equipped with the supremum
norm. The contraction mapping theorem (Banach fixed-point theorem, which we prove
in the next lesson) guarantees that $T$ has a unique fixed point and that
$\|T^n V_0 - V^*\|_\infty \le \gamma^n \|V_0 - V^*\|_\infty$. Understanding
\emph{why} this works requires every concept in this lesson: metric, norm,
completeness, and the operator framework.
\end{rlconnection}

%% ============================================================
%% SECTION 2: METRIC SPACES
%% ============================================================
\section{Metric Spaces}
\label{sec:metric}

\subsection{Definition and Examples}
\label{subsec:metric_def}

\begin{definition}[Metric Space]
\label{def:metric}
A \emph{metric} on a set $X$ is a function $d \colon X \times X \to [0,\infty)$
satisfying, for all $x, y, z \in X$:
\begin{enumerate}[nosep,label=(M\arabic*)]
    \item \label{M1} $d(x,y) = 0$ if and only if $x = y$.
    \hfill(identity of indiscernibles)
    \item \label{M2} $d(x,y) = d(y,x)$.
    \hfill(symmetry)
    \item \label{M3} $d(x,z) \le d(x,y) + d(y,z)$.
    \hfill(triangle inequality)
\end{enumerate}
The pair $(X,d)$ is called a \emph{metric space}. Elements of $X$ are called
\emph{points}.
\end{definition}

\begin{remark}
Non-negativity $d(x,y) \ge 0$ follows from the axioms: by \ref{M3} and \ref{M1},
$0 = d(x,x) \le d(x,y) + d(y,x) = 2d(x,y)$ by \ref{M2}, so $d(x,y) \ge 0$.
\end{remark}

\begin{example}[Euclidean Metric]
\label{ex:euclidean}
$X = \R^n$ with $d(x,y) = \bigl(\sum_{i=1}^n (x_i - y_i)^2\bigr)^{1/2}$. The
triangle inequality follows from the Cauchy--Schwarz inequality.
\end{example}

\begin{example}[$p$-Metrics on $\R^n$]
\label{ex:p_metric}
For $1 \le p \le \infty$, define $d_p(x,y) = \|x - y\|_p$ where
\[
\|x\|_p =
\begin{cases}
\bigl(\sum_{i=1}^n |x_i|^p\bigr)^{1/p} & \text{if } 1 \le p < \infty,\\
\max_{1 \le i \le n} |x_i| & \text{if } p = \infty.
\end{cases}
\]
Each $d_p$ is a metric. The triangle inequality for $1 < p < \infty$ is
Minkowski's inequality (proved in Proposition~\ref{prop:minkowski}).
\end{example}

\begin{example}[Discrete Metric]
\label{ex:discrete}
On any set $X$, define $d(x,y) = 0$ if $x = y$ and $d(x,y) = 1$ if $x \ne y$.
This is a metric. Every subset of $X$ is open (and closed), so $X$ with the
discrete metric is a \emph{discrete topological space}.
\end{example}

\begin{example}[Function Space $C[a,b]$]
\label{ex:Cab}
Let $C[a,b]$ denote the set of continuous real-valued functions on $[a,b]$. Define
\[
d_\infty(f,g) = \max_{t \in [a,b]} |f(t) - g(t)|, \qquad
d_1(f,g) = \int_a^b |f(t) - g(t)|\, dt.
\]
Both are metrics on $C[a,b]$. The metric $d_\infty$ measures uniform distance;
$d_1$ measures average distance.
\end{example}

\begin{example}[Sequence Spaces $\ell^p$]
\label{ex:ell_p}
For $1 \le p < \infty$, define $\ell^p = \{(x_n)_{n=1}^\infty \subset \R :
\sum_{n=1}^\infty |x_n|^p < \infty\}$ with metric
$d(x,y) = \bigl(\sum_{n=1}^\infty |x_n - y_n|^p\bigr)^{1/p}$.
For $p = \infty$, $\ell^\infty = \{(x_n) : \sup_n |x_n| < \infty\}$ with
$d(x,y) = \sup_n |x_n - y_n|$.
\end{example}

\begin{rlconnection}[State Spaces as Metric Spaces]
In RL with a finite state space $\mathcal{S} = \{s_1, \ldots, s_n\}$, a value
function $V \colon \mathcal{S} \to \R$ is simply a vector in $\R^n$. The space of
all value functions is $(\R^n, \|\cdot\|_\infty)$, which is a metric space. For
infinite or continuous state spaces, value functions live in function spaces like
$C(\mathcal{S})$ or $\ell^\infty(\mathcal{S})$, and the theory of this section
becomes essential.
\end{rlconnection}

\subsection{Topology of Metric Spaces}
\label{subsec:topology}

\begin{definition}[Open Ball]
\label{def:open_ball}
In a metric space $(X,d)$, the \emph{open ball} of radius $r > 0$ centered at
$x \in X$ is
\[
B(x,r) = \{y \in X : d(x,y) < r\}.
\]
The \emph{closed ball} is $\overline{B}(x,r) = \{y \in X : d(x,y) \le r\}$.
\end{definition}

\begin{definition}[Open and Closed Sets]
\label{def:open_closed}
Let $(X,d)$ be a metric space.
\begin{enumerate}[nosep,label=(\roman*)]
    \item A set $G \subseteq X$ is \emph{open} if for every $x \in G$ there exists
    $r > 0$ such that $B(x,r) \subseteq G$.
    \item A set $F \subseteq X$ is \emph{closed} if its complement $X \setminus F$
    is open.
\end{enumerate}
\end{definition}

\begin{proposition}[Properties of Open and Closed Sets]
\label{prop:open_closed}
In a metric space $(X,d)$:
\begin{enumerate}[nosep,label=(\alph*)]
    \item $\emptyset$ and $X$ are both open and closed.
    \item Any union of open sets is open.
    \item Any finite intersection of open sets is open.
    \item Any intersection of closed sets is closed.
    \item Any finite union of closed sets is closed.
    \item Every open ball $B(x,r)$ is open.
\end{enumerate}
\end{proposition}

\begin{proof}
\textbf{(a)} $\emptyset$ is open vacuously. $X$ is open because $B(x,r) \subseteq X$
for all $x$ and $r > 0$.

\textbf{(b)} Let $\{G_\alpha\}_{\alpha \in I}$ be a family of open sets and let
$x \in \bigcup_\alpha G_\alpha$. Then $x \in G_\beta$ for some $\beta$, so there
exists $r > 0$ with $B(x,r) \subseteq G_\beta \subseteq \bigcup_\alpha G_\alpha$.

\textbf{(c)} Let $G_1, \ldots, G_n$ be open and let $x \in \bigcap_{i=1}^n G_i$.
For each $i$, there exists $r_i > 0$ with $B(x,r_i) \subseteq G_i$. Set
$r = \min(r_1, \ldots, r_n) > 0$. Then $B(x,r) \subseteq \bigcap_{i=1}^n G_i$.

\textbf{(d)--(e)} Follow from (b)--(c) by taking complements (De Morgan's laws).

\textbf{(f)} Let $y \in B(x,r)$ and set $\varepsilon = r - d(x,y) > 0$. For any
$z \in B(y,\varepsilon)$, we have $d(x,z) \le d(x,y) + d(y,z) < d(x,y) +
\varepsilon = r$, so $z \in B(x,r)$. Thus $B(y,\varepsilon) \subseteq B(x,r)$.
\end{proof}

\begin{definition}[Interior, Closure, Boundary]
\label{def:int_cl_bd}
Let $A \subseteq X$ in a metric space $(X,d)$.
\begin{enumerate}[nosep,label=(\roman*)]
    \item The \emph{interior} of $A$ is $\Int(A) = \bigcup \{G \subseteq A : G \text{ is open}\}$.
    \item The \emph{closure} of $A$ is $\cl(A) = \bigcap \{F \supseteq A : F \text{ is closed}\}$.
    \item The \emph{boundary} of $A$ is $\partial A = \cl(A) \setminus \Int(A)$.
    \item $A$ is \emph{dense} in $X$ if $\cl(A) = X$.
\end{enumerate}
\end{definition}

\begin{proposition}[Closure Characterization]
\label{prop:closure_char}
$x \in \cl(A)$ if and only if for every $r > 0$, $B(x,r) \cap A \ne \emptyset$.
Equivalently, $x \in \cl(A)$ if and only if there exists a sequence $(x_n)$ in $A$
with $x_n \to x$.
\end{proposition}

\begin{proof}
($\Rightarrow$) Suppose $x \in \cl(A)$ and there exists $r > 0$ with
$B(x,r) \cap A = \emptyset$. Then $A \subseteq X \setminus B(x,r)$, which is
closed. So $\cl(A) \subseteq X \setminus B(x,r)$, contradicting $x \in \cl(A)$.

($\Leftarrow$) Suppose every ball around $x$ meets $A$, and let $F$ be any closed
set containing $A$. If $x \notin F$, then $x \in X \setminus F$, which is open,
so there exists $r > 0$ with $B(x,r) \subseteq X \setminus F$. But then
$B(x,r) \cap A \subseteq B(x,r) \cap F = \emptyset$, contradicting our assumption.
So $x \in F$ for every closed $F \supseteq A$, hence $x \in \cl(A)$.

For the sequential characterization: if $B(x,1/n) \cap A \ne \emptyset$ for each
$n$, pick $x_n \in B(x,1/n) \cap A$; then $d(x_n,x) < 1/n \to 0$. Conversely, if
$x_n \to x$ with each $x_n \in A$, then every ball $B(x,r)$ contains $x_n$ for
large $n$, so $B(x,r) \cap A \ne \emptyset$.
\end{proof}

\subsection{Convergence and Completeness}
\label{subsec:convergence}

\begin{definition}[Convergence]
\label{def:convergence}
A sequence $(x_n)$ in a metric space $(X,d)$ \emph{converges} to $x \in X$ if
$d(x_n, x) \to 0$ as $n \to \infty$, i.e., for every $\varepsilon > 0$ there
exists $N \in \N$ such that $d(x_n, x) < \varepsilon$ for all $n \ge N$. We write
$x_n \to x$ or $\lim_{n \to \infty} x_n = x$.
\end{definition}

\begin{proposition}[Uniqueness of Limits]
\label{prop:unique_limit}
In a metric space, limits are unique: if $x_n \to x$ and $x_n \to y$, then $x = y$.
\end{proposition}

\begin{proof}
$0 \le d(x,y) \le d(x,x_n) + d(x_n,y) \to 0 + 0 = 0$, so $d(x,y) = 0$, hence $x = y$.
\end{proof}

\begin{definition}[Cauchy Sequence]
\label{def:cauchy}
A sequence $(x_n)$ in $(X,d)$ is a \emph{Cauchy sequence} if for every
$\varepsilon > 0$ there exists $N \in \N$ such that $d(x_m, x_n) < \varepsilon$
for all $m, n \ge N$.
\end{definition}

\begin{proposition}
\label{prop:convergent_is_cauchy}
Every convergent sequence is Cauchy.
\end{proposition}

\begin{proof}
If $x_n \to x$, then for $\varepsilon > 0$ there exists $N$ such that
$d(x_n, x) < \varepsilon/2$ for $n \ge N$. For $m, n \ge N$:
$d(x_m, x_n) \le d(x_m, x) + d(x, x_n) < \varepsilon/2 + \varepsilon/2 = \varepsilon$.
\end{proof}

\begin{warning}[Cauchy Does Not Imply Convergent]
The converse fails in general. In $\Q$ with the usual metric, the sequence of
rational approximations to $\sqrt{2}$ (e.g., defined by $x_1 = 1$,
$x_{n+1} = (x_n + 2/x_n)/2$) is Cauchy but does not converge in $\Q$.
The ``missing limit'' is $\sqrt{2} \notin \Q$.
\end{warning}

\begin{definition}[Complete Metric Space]
\label{def:complete}
A metric space $(X,d)$ is \emph{complete} if every Cauchy sequence in $X$
converges to a limit in $X$.
\end{definition}

\begin{theorem}[Completeness of $\R$]
\label{thm:R_complete}
$(\R, |\cdot|)$ is a complete metric space.
\end{theorem}

\begin{proof}
Let $(x_n)$ be a Cauchy sequence in $\R$. First, $(x_n)$ is bounded: there exists
$N$ such that $|x_m - x_n| < 1$ for all $m, n \ge N$, so
$|x_n| \le |x_N| + 1$ for $n \ge N$, giving
$|x_n| \le \max(|x_1|, \ldots, |x_{N-1}|, |x_N| + 1)$ for all $n$.

By the Bolzano--Weierstrass theorem (every bounded sequence in $\R$ has a
convergent subsequence), there is a subsequence $x_{n_k} \to x$ for some $x \in \R$.

We show $x_n \to x$. Given $\varepsilon > 0$, choose $N_1$ such that
$|x_m - x_n| < \varepsilon/2$ for $m, n \ge N_1$, and $N_2$ such that
$|x_{n_k} - x| < \varepsilon/2$ for $k \ge N_2$. For $n \ge N_1$, pick $k$ with
$n_k \ge \max(N_1, n_{N_2})$. Then
$|x_n - x| \le |x_n - x_{n_k}| + |x_{n_k} - x| < \varepsilon/2 + \varepsilon/2 = \varepsilon$.
\end{proof}

\begin{theorem}[Completeness of $\R^n$]
\label{thm:Rn_complete}
$(\R^n, d_p)$ is complete for every $1 \le p \le \infty$.
\end{theorem}

\begin{proof}
Let $(x^{(k)})_{k=1}^\infty$ be a Cauchy sequence in $(\R^n, d_p)$, where
$x^{(k)} = (x_1^{(k)}, \ldots, x_n^{(k)})$. For each coordinate $i$,
$|x_i^{(k)} - x_i^{(m)}| \le \|x^{(k)} - x^{(m)}\|_p$ (since the $\ell^p$ norm
dominates each coordinate), so $(x_i^{(k)})_{k=1}^\infty$ is Cauchy in $\R$.
By Theorem~\ref{thm:R_complete}, $x_i^{(k)} \to x_i$ for some $x_i \in \R$.
Set $x = (x_1, \ldots, x_n)$. Then
$\|x^{(k)} - x\|_p = \bigl(\sum_{i=1}^n |x_i^{(k)} - x_i|^p\bigr)^{1/p} \to 0$
as $k \to \infty$ (each term goes to $0$; this is a finite sum). So $x^{(k)} \to x$.
\end{proof}

\begin{definition}[Dense Subset, Separable Space]
\label{def:separable}
A metric space $(X,d)$ is \emph{separable} if it contains a countable dense subset,
i.e., there exists a countable set $D \subseteq X$ with $\cl(D) = X$.
\end{definition}

\begin{example}
$\R^n$ is separable: $\Q^n$ is countable and dense. The space $C[a,b]$ with the
sup metric is separable: by the Weierstrass approximation theorem, polynomials with
rational coefficients form a countable dense subset.
\end{example}

\begin{theorem}[Completion of a Metric Space]
\label{thm:completion}
Every metric space $(X,d)$ has a \emph{completion}: there exists a complete metric
space $(\hat{X}, \hat{d})$ and an isometry $\iota \colon X \to \hat{X}$ (i.e.,
$\hat{d}(\iota(x), \iota(y)) = d(x,y)$ for all $x, y \in X$) such that
$\iota(X)$ is dense in $\hat{X}$. The completion is unique up to isometry.
\end{theorem}

\begin{proof}
\textbf{Construction.}
Let $\mathcal{C}$ be the set of all Cauchy sequences in $X$. Define an equivalence
relation: $(x_n) \sim (y_n)$ if $d(x_n, y_n) \to 0$. Let
$\hat{X} = \mathcal{C}/{\sim}$ be the set of equivalence classes, and write
$[(x_n)]$ for the class of $(x_n)$.

\textbf{The metric is well defined.}
For $(x_n), (y_n) \in \mathcal{C}$, the sequence $(d(x_n, y_n))$ is Cauchy in $\R$
(since $|d(x_m,y_m) - d(x_n,y_n)| \le d(x_m,x_n) + d(y_m,y_n)$ by the triangle
inequality), hence convergent. Define
$\hat{d}([(x_n)], [(y_n)]) = \lim_{n \to \infty} d(x_n, y_n)$.
One checks that this is independent of representatives (if $(x_n) \sim (x_n')$
and $(y_n) \sim (y_n')$, then $|d(x_n,y_n) - d(x_n',y_n')| \le d(x_n,x_n') +
d(y_n,y_n') \to 0$) and satisfies the metric axioms.

\textbf{Isometric embedding.}
Define $\iota(x) = [(x, x, x, \ldots)]$ (the class of the constant sequence).
Then $\hat{d}(\iota(x), \iota(y)) = \lim_n d(x,y) = d(x,y)$, so $\iota$ is an
isometry.

\textbf{Density of $\iota(X)$.}
Given $[(x_n)] \in \hat{X}$ and $\varepsilon > 0$, choose $N$ with
$d(x_m, x_n) < \varepsilon/2$ for $m, n \ge N$. Then
$\hat{d}([(x_n)], \iota(x_N)) = \lim_n d(x_n, x_N) \le \varepsilon/2 < \varepsilon$.

\textbf{Completeness of $\hat{X}$.}
Let $(p_k)$ be Cauchy in $\hat{X}$. Since $\iota(X)$ is dense, for each $k$ choose
$z_k \in X$ with $\hat{d}(p_k, \iota(z_k)) < 1/k$. We claim $(z_k)$ is Cauchy in
$X$: indeed, $d(z_k, z_m) = \hat{d}(\iota(z_k), \iota(z_m)) \le
\hat{d}(\iota(z_k), p_k) + \hat{d}(p_k, p_m) + \hat{d}(p_m, \iota(z_m)) <
1/k + \hat{d}(p_k, p_m) + 1/m \to 0$. Let $p = [(z_k)] \in \hat{X}$. Then
$\hat{d}(p_k, p) \le \hat{d}(p_k, \iota(z_k)) + \hat{d}(\iota(z_k), p) <
1/k + \hat{d}(\iota(z_k), p) \to 0$ since
$\hat{d}(\iota(z_k), p) = \lim_n d(z_k, z_n) \to 0$ as $k \to \infty$.

\textbf{Uniqueness.}
If $(\hat{X}_1, \hat{d}_1, \iota_1)$ and $(\hat{X}_2, \hat{d}_2, \iota_2)$ are
two completions, define $\varphi \colon \iota_1(X) \to \iota_2(X)$ by
$\varphi(\iota_1(x)) = \iota_2(x)$. This is an isometry, hence uniformly
continuous, and extends uniquely to an isometry $\hat{X}_1 \to \hat{X}_2$ by
density and completeness.
\end{proof}

\begin{rlconnection}[Completion and Function Approximation]
In RL with function approximation, we often work with a parametric class of
functions (e.g., neural networks) that is not complete under the sup norm. The
completion theorem guarantees that this class can be embedded into a complete
space. Understanding which space we ``complete to'' is crucial for analyzing
whether approximate value iteration converges.
\end{rlconnection}

%% ============================================================
%% SECTION 3: NORMED SPACES
%% ============================================================
\section{Normed Spaces}
\label{sec:normed}

\subsection{Definition and Basic Properties}
\label{subsec:norm_def}

\begin{definition}[Normed Space]
\label{def:normed}
A \emph{normed space} (or \emph{normed vector space}) is a pair $(X, \|\cdot\|)$
where $X$ is a vector space over $\R$ (or $\C$) and
$\|\cdot\| \colon X \to [0,\infty)$ is a \emph{norm}, i.e., for all
$x, y \in X$ and $\alpha$ scalar:
\begin{enumerate}[nosep,label=(N\arabic*)]
    \item \label{N1} $\|x\| = 0$ if and only if $x = 0$.
    \hfill(positive definiteness)
    \item \label{N2} $\|\alpha x\| = |\alpha| \, \|x\|$.
    \hfill(absolute homogeneity)
    \item \label{N3} $\|x + y\| \le \|x\| + \|y\|$.
    \hfill(triangle inequality)
\end{enumerate}
\end{definition}

\begin{proposition}[Every Norm Induces a Metric]
\label{prop:norm_metric}
If $(X, \|\cdot\|)$ is a normed space, then $d(x,y) = \|x - y\|$ defines a metric
on $X$. This metric is \emph{translation-invariant} ($d(x+z, y+z) = d(x,y)$) and
\emph{homogeneous} ($d(\alpha x, \alpha y) = |\alpha| d(x,y)$).
\end{proposition}

\begin{proof}
\ref{M1}: $d(x,y) = \|x - y\| = 0 \iff x - y = 0 \iff x = y$.
\ref{M2}: $d(x,y) = \|x - y\| = \|{-(y-x)}\| = |{-1}| \cdot \|y - x\| = d(y,x)$.
\ref{M3}: $d(x,z) = \|x - z\| = \|(x-y)+(y-z)\| \le \|x-y\| + \|y-z\| = d(x,y) + d(y,z)$.
Translation invariance: $d(x+z, y+z) = \|(x+z)-(y+z)\| = \|x-y\| = d(x,y)$.
Homogeneity: $d(\alpha x, \alpha y) = \|\alpha x - \alpha y\| = |\alpha| \|x - y\| = |\alpha| d(x,y)$.
\end{proof}

\begin{remark}
Not every metric comes from a norm. The discrete metric on $\R^n$ (with $n \ge 1$)
is not translation-invariant (take $x = 0$, $y = z$ with $z \ne 0$:
$d(x,y) = 1$ but $d(x+z, y+z) = d(z, 2z) = 1$, which happens to be invariant in
this case, but it is not homogeneous: $d(2x, 2y) = 1 \ne 2 = 2 d(x,y)$ when
$x \ne y$). A metric arises from a norm if and only if it is translation-invariant
and absolutely homogeneous.
\end{remark}

\subsection{H\"{o}lder's and Minkowski's Inequalities}

To show that $\ell^p$ norms are indeed norms, we need the following fundamental
inequalities.

\begin{lemma}[Young's Inequality]
\label{lem:young}
For $a, b \ge 0$ and $p, q > 1$ with $\frac{1}{p} + \frac{1}{q} = 1$:
\[
ab \le \frac{a^p}{p} + \frac{b^q}{q}.
\]
\end{lemma}

\begin{proof}
The function $\varphi(t) = e^t$ is convex. By the definition of convexity,
$e^{\lambda s + (1-\lambda)t} \le \lambda e^s + (1-\lambda) e^t$ for
$\lambda \in [0,1]$. Set $\lambda = 1/p$, $s = p \ln a$, $t = q \ln b$ (assuming
$a, b > 0$):
\[
ab = e^{\ln a + \ln b} = e^{(1/p)(p \ln a) + (1/q)(q \ln b)} \le
\frac{1}{p} e^{p \ln a} + \frac{1}{q} e^{q \ln b} = \frac{a^p}{p} + \frac{b^q}{q}.
\]
The case $a = 0$ or $b = 0$ is trivial.
\end{proof}

\begin{proposition}[H\"{o}lder's Inequality]
\label{prop:holder}
Let $1 < p < \infty$ and $q = p/(p-1)$ (so $1/p + 1/q = 1$). For
$x = (x_i), y = (y_i) \in \R^n$ (or sequences in $\ell^p, \ell^q$):
\[
\sum_{i} |x_i y_i| \le \|x\|_p \, \|y\|_q.
\]
\end{proposition}

\begin{proof}
If $\|x\|_p = 0$ or $\|y\|_q = 0$, the result is trivial. Otherwise, set
$a_i = |x_i|/\|x\|_p$ and $b_i = |y_i|/\|y\|_q$. By Young's inequality:
$a_i b_i \le a_i^p/p + b_i^q/q$. Summing:
\[
\sum_i a_i b_i \le \frac{1}{p} \sum_i \frac{|x_i|^p}{\|x\|_p^p} +
\frac{1}{q} \sum_i \frac{|y_i|^q}{\|y\|_q^q} = \frac{1}{p} + \frac{1}{q} = 1.
\]
Multiplying both sides by $\|x\|_p \|y\|_q$ gives the result.
\end{proof}

\begin{proposition}[Minkowski's Inequality]
\label{prop:minkowski}
For $1 \le p < \infty$ and $x, y \in \R^n$ (or $\ell^p$):
$\|x + y\|_p \le \|x\|_p + \|y\|_p$.
\end{proposition}

\begin{proof}
The case $p = 1$ is the ordinary triangle inequality for absolute values. For
$p > 1$, let $q = p/(p-1)$. We have:
\begin{align*}
\|x + y\|_p^p &= \sum_i |x_i + y_i|^p = \sum_i |x_i + y_i|^{p-1} |x_i + y_i| \\
&\le \sum_i |x_i + y_i|^{p-1} |x_i| + \sum_i |x_i + y_i|^{p-1} |y_i|.
\end{align*}
Apply H\"{o}lder's inequality to each sum with exponents $q$ and $p$. Note that
$(p-1)q = p$, so $\bigl(\sum_i |x_i + y_i|^{(p-1)q}\bigr)^{1/q} =
\|x + y\|_p^{p/q}$. Therefore:
\[
\|x + y\|_p^p \le \|x + y\|_p^{p/q} \bigl(\|x\|_p + \|y\|_p\bigr).
\]
Dividing both sides by $\|x + y\|_p^{p/q}$ (if nonzero) and using $p - p/q = 1$
gives $\|x + y\|_p \le \|x\|_p + \|y\|_p$.
\end{proof}

\subsection{Important Examples of Normed Spaces}
\label{subsec:norm_examples}

\begin{example}[$\R^n$ with $p$-norms]
\label{ex:Rn_norms}
For $1 \le p \le \infty$, $(\R^n, \|\cdot\|_p)$ is a normed space.
Positive definiteness and homogeneity are immediate; the triangle inequality
for $1 \le p < \infty$ is Minkowski's inequality. For $p = \infty$,
$\|x + y\|_\infty = \max_i |x_i + y_i| \le \max_i (|x_i| + |y_i|) \le
\max_i |x_i| + \max_i |y_i| = \|x\|_\infty + \|y\|_\infty$.
\end{example}

\begin{example}[$C[a,b]$ with Supremum Norm]
\label{ex:Cab_sup}
$C[a,b]$ with $\|f\|_\infty = \max_{t \in [a,b]} |f(t)|$ is a normed space. The
maximum exists because $f$ is continuous on the compact set $[a,b]$.
\end{example}

\begin{example}[$C[a,b]$ with $L^1$ Norm]
\label{ex:Cab_L1}
$C[a,b]$ with $\|f\|_1 = \int_a^b |f(t)|\, dt$ is a normed space. This norm
captures ``average'' rather than ``worst-case'' behavior.
\end{example}

\begin{example}[Sequence Spaces $\ell^p$]
\label{ex:ell_p_norm}
For $1 \le p \le \infty$, $(\ell^p, \|\cdot\|_p)$ is a normed space by
Minkowski's inequality (extended to infinite sums).
\end{example}

\subsection{Equivalence of Norms in Finite Dimensions}
\label{subsec:norm_equiv}

\begin{definition}[Equivalent Norms]
\label{def:norm_equiv}
Two norms $\|\cdot\|_a$ and $\|\cdot\|_b$ on a vector space $X$ are
\emph{equivalent} if there exist constants $c, C > 0$ such that
\[
c \|x\|_a \le \|x\|_b \le C \|x\|_a \quad \text{for all } x \in X.
\]
\end{definition}

\begin{remark}
Equivalent norms define the same topology: they determine the same open sets,
convergent sequences, Cauchy sequences, and continuous functions. In particular,
a space is complete under one norm if and only if it is complete under any
equivalent norm.
\end{remark}

\begin{theorem}[All Norms on $\R^n$ Are Equivalent]
\label{thm:norm_equiv}
On $\R^n$, every two norms are equivalent.
\end{theorem}

\begin{proof}
It suffices to show that any norm $\|\cdot\|$ on $\R^n$ is equivalent to
$\|\cdot\|_\infty$, since equivalence is transitive.

\textbf{Step 1: $\|\cdot\|$ is bounded above by a multiple of $\|\cdot\|_\infty$.}
Let $e_1, \ldots, e_n$ be the standard basis. For $x = \sum_{i=1}^n x_i e_i$:
\[
\|x\| = \Bigl\|\sum_{i=1}^n x_i e_i\Bigr\| \le \sum_{i=1}^n |x_i| \, \|e_i\|
\le \|x\|_\infty \sum_{i=1}^n \|e_i\| = C \|x\|_\infty,
\]
where $C = \sum_{i=1}^n \|e_i\|$.

\textbf{Step 2: $\|\cdot\|$ is continuous with respect to $\|\cdot\|_\infty$.}
By Step 1, $|\|x\| - \|y\|| \le \|x - y\| \le C \|x - y\|_\infty$, so
$x \mapsto \|x\|$ is Lipschitz (hence continuous) on $(\R^n, \|\cdot\|_\infty)$.

\textbf{Step 3: $\|\cdot\|$ is bounded below by a multiple of $\|\cdot\|_\infty$.}
The unit sphere $S = \{x \in \R^n : \|x\|_\infty = 1\}$ is closed and bounded in
$(\R^n, \|\cdot\|_\infty)$, hence compact (by the Heine--Borel theorem). The
continuous function $x \mapsto \|x\|$ attains its minimum on $S$: let
$c = \min_{x \in S} \|x\|$. Since $0 \notin S$ (because $\|x\|_\infty = 1$ implies
$x \ne 0$, so $\|x\| > 0$ by \ref{N1}), we have $c > 0$. For any $x \ne 0$,
$x/\|x\|_\infty \in S$, so $\|x/\|x\|_\infty\| \ge c$, giving
$\|x\| \ge c \|x\|_\infty$.

Combining Steps 1 and 3: $c \|x\|_\infty \le \|x\| \le C \|x\|_\infty$.
\end{proof}

\begin{keyresult}[Finite-Dimensional Norm Equivalence]
On a finite-dimensional vector space, all norms are equivalent. This means
convergence, completeness, compactness, and continuity are norm-independent in
$\R^n$. This fails dramatically in infinite dimensions (see
Warning~\ref{warn:inf_dim_norms}).
\end{keyresult}

\begin{corollary}
\label{cor:finite_dim_complete}
Every finite-dimensional normed space is complete (i.e., is a Banach space).
\end{corollary}

\begin{proof}
Any norm on $\R^n$ is equivalent to $\|\cdot\|_\infty$, and $(\R^n, \|\cdot\|_\infty)$
is complete by Theorem~\ref{thm:Rn_complete}. Since equivalent norms have the
same Cauchy sequences and convergent sequences, completeness is preserved.
\end{proof}

\begin{warning}[Infinite-Dimensional Spaces Are Different]
\label{warn:inf_dim_norms}
On $C[0,1]$, the norms $\|f\|_\infty = \max |f|$ and $\|f\|_1 = \int_0^1 |f|$
are \emph{not} equivalent. Consider $f_n(t) = t^n$: then $\|f_n\|_\infty = 1$ but
$\|f_n\|_1 = 1/(n+1) \to 0$. If the norms were equivalent, $\|f_n\|_1 \to 0$
would imply $\|f_n\|_\infty \to 0$, a contradiction. This has profound consequences:
$C[0,1]$ is complete under $\|\cdot\|_\infty$ but \emph{not} under $\|\cdot\|_1$.
\end{warning}

\begin{rlconnection}[Norm Choice in RL]
The choice of norm matters for the convergence properties of RL algorithms.
Value iteration contracts in the $\|\cdot\|_\infty$ norm with factor $\gamma$,
but the same operator need not contract in the $\|\cdot\|_2$ norm. For finite
state spaces all norms are equivalent (Theorem~\ref{thm:norm_equiv}), so
convergence in one norm implies convergence in all. For infinite state spaces
with function approximation, the norm matters critically, and
non-equivalence of norms (Warning~\ref{warn:inf_dim_norms}) is the root of
many difficulties in approximate dynamic programming.
\end{rlconnection}

%% ============================================================
%% SECTION 4: BANACH SPACES
%% ============================================================
\section{Banach Spaces}
\label{sec:banach}

\begin{definition}[Banach Space]
\label{def:banach}
A \emph{Banach space} is a complete normed space, i.e., a normed space $(X, \|\cdot\|)$
that is complete as a metric space under the induced metric $d(x,y) = \|x - y\|$.
\end{definition}

\subsection{Examples and Non-Examples}

\begin{theorem}[$\ell^p$ is a Banach Space]
\label{thm:ell_p_banach}
For $1 \le p \le \infty$, $(\ell^p, \|\cdot\|_p)$ is a Banach space.
\end{theorem}

\begin{proof}
We prove the case $1 \le p < \infty$; the case $p = \infty$ is analogous.
Let $(x^{(k)})_{k=1}^\infty$ be a Cauchy sequence in $\ell^p$, where
$x^{(k)} = (x_1^{(k)}, x_2^{(k)}, \ldots)$.

\textbf{Step 1: Identify the candidate limit.}
For each fixed $i$, $|x_i^{(k)} - x_i^{(m)}| \le \|x^{(k)} - x^{(m)}\|_p \to 0$,
so $(x_i^{(k)})_k$ is Cauchy in $\R$, hence convergent. Define
$x_i = \lim_{k \to \infty} x_i^{(k)}$ and $x = (x_1, x_2, \ldots)$.

\textbf{Step 2: Show $x \in \ell^p$.}
Fix $\varepsilon > 0$ and choose $K$ such that $\|x^{(k)} - x^{(m)}\|_p < \varepsilon$
for $k, m \ge K$. For any finite $N$:
\[
\Bigl(\sum_{i=1}^N |x_i^{(k)} - x_i^{(m)}|^p\Bigr)^{1/p} \le
\|x^{(k)} - x^{(m)}\|_p < \varepsilon.
\]
Taking $m \to \infty$ (with $N$ fixed): $\bigl(\sum_{i=1}^N |x_i^{(k)} - x_i|^p\bigr)^{1/p}
\le \varepsilon$. Since this holds for all $N$,
$\|x^{(k)} - x\|_p \le \varepsilon < \infty$. By the triangle inequality,
$\|x\|_p \le \|x - x^{(k)}\|_p + \|x^{(k)}\|_p < \infty$, so $x \in \ell^p$.

\textbf{Step 3: Show $x^{(k)} \to x$ in $\ell^p$.}
From Step 2, $\|x^{(k)} - x\|_p \le \varepsilon$ for all $k \ge K$, so
$x^{(k)} \to x$ in $\ell^p$.
\end{proof}

\begin{theorem}[$C[a,b]$ is Banach Under the Sup Norm]
\label{thm:Cab_banach}
$(C[a,b], \|\cdot\|_\infty)$ is a Banach space.
\end{theorem}

\begin{proof}
Let $(f_n)$ be Cauchy in $(C[a,b], \|\cdot\|_\infty)$. For each fixed $t \in [a,b]$,
$|f_n(t) - f_m(t)| \le \|f_n - f_m\|_\infty \to 0$, so $(f_n(t))$ is Cauchy in
$\R$, hence convergent. Define $f(t) = \lim_{n \to \infty} f_n(t)$.

\textbf{Uniform convergence.}
Given $\varepsilon > 0$, choose $N$ such that $\|f_n - f_m\|_\infty < \varepsilon$
for $n, m \ge N$. For each $t$, $|f_n(t) - f_m(t)| < \varepsilon$. Letting
$m \to \infty$: $|f_n(t) - f(t)| \le \varepsilon$ for all $t$ and $n \ge N$.
Hence $\|f_n - f\|_\infty \le \varepsilon$, so $f_n \to f$ uniformly.

\textbf{Continuity of $f$.}
The uniform limit of continuous functions is continuous: for any $t_0 \in [a,b]$ and
$\varepsilon > 0$, choose $n$ with $\|f_n - f\|_\infty < \varepsilon/3$, then
$\delta > 0$ with $|f_n(t) - f_n(t_0)| < \varepsilon/3$ for $|t - t_0| < \delta$.
Then $|f(t) - f(t_0)| \le |f(t) - f_n(t)| + |f_n(t) - f_n(t_0)| +
|f_n(t_0) - f(t_0)| < \varepsilon$.

So $f \in C[a,b]$ and $f_n \to f$ in $\|\cdot\|_\infty$.
\end{proof}

\begin{example}[$C[a,b]$ Is Not Banach Under $\|\cdot\|_1$]
\label{ex:Cab_not_banach}
The space $(C[0,1], \|\cdot\|_1)$ is \emph{not} complete. Define
$f_n \colon [0,1] \to \R$ by
\[
f_n(t) = \begin{cases}
0 & \text{if } 0 \le t \le 1/2,\\
n(t - 1/2) & \text{if } 1/2 < t \le 1/2 + 1/n,\\
1 & \text{if } 1/2 + 1/n < t \le 1.
\end{cases}
\]
Each $f_n$ is continuous, and $(f_n)$ is Cauchy under $\|\cdot\|_1$ (since
$\|f_n - f_m\|_1 \le 1/\min(m,n) \to 0$). But the pointwise limit is the
discontinuous step function $\indicator_{(1/2, 1]}$, which is not in $C[0,1]$.
No continuous function can be the $\|\cdot\|_1$-limit either, since $\|f_n - g\|_1
\to 0$ and $\|f_n - \indicator_{(1/2,1]}\|_1 \to 0$ would give
$g = \indicator_{(1/2,1]}$ a.e., contradicting continuity of $g$.
\end{example}

\subsection{Series in Banach Spaces}

\begin{definition}[Absolute Convergence]
\label{def:abs_conv}
A series $\sum_{n=1}^\infty x_n$ in a normed space is \emph{absolutely convergent}
if $\sum_{n=1}^\infty \|x_n\| < \infty$. It is \emph{(unconditionally) convergent}
if the partial sums $S_N = \sum_{n=1}^N x_n$ converge in norm.
\end{definition}

\begin{theorem}[Characterization of Banach Spaces via Series]
\label{thm:banach_series}
A normed space $X$ is a Banach space if and only if every absolutely convergent
series in $X$ converges.
\end{theorem}

\begin{proof}
($\Rightarrow$) Suppose $X$ is Banach and $\sum \|x_n\| < \infty$. The partial
sums $S_N = \sum_{n=1}^N x_n$ form a Cauchy sequence: for $M > N$,
\[
\|S_M - S_N\| = \Bigl\|\sum_{n=N+1}^M x_n\Bigr\| \le \sum_{n=N+1}^M \|x_n\| \to 0
\]
as $N, M \to \infty$ (tail of a convergent series). Since $X$ is complete, $(S_N)$
converges.

($\Leftarrow$) Suppose every absolutely convergent series converges. Let $(y_n)$
be a Cauchy sequence. Choose a subsequence $(y_{n_k})$ with
$\|y_{n_{k+1}} - y_{n_k}\| < 2^{-k}$ (possible since $(y_n)$ is Cauchy). Define
$x_1 = y_{n_1}$ and $x_k = y_{n_k} - y_{n_{k-1}}$ for $k \ge 2$. Then the partial
sums satisfy $\sum_{k=1}^K x_k = y_{n_K}$, and
$\sum_{k=2}^\infty \|x_k\| \le \sum_{k=1}^\infty 2^{-k} = 1 < \infty$. So
$\sum x_k$ converges absolutely, hence converges by hypothesis: $y_{n_K} \to y$ for
some $y \in X$. Since $(y_n)$ is Cauchy and has a convergent subsequence, the full
sequence converges to $y$ (by the argument: $\|y_n - y\| \le
\|y_n - y_{n_K}\| + \|y_{n_K} - y\| \to 0$ for suitable $K$).
\end{proof}

\begin{keyresult}[Banach Space Characterization]
A normed space is Banach if and only if every absolutely convergent series
converges. This characterization is often the most convenient way to verify
completeness: one constructs a convergent series rather than chasing arbitrary
Cauchy sequences.
\end{keyresult}

\begin{rlconnection}[Why Completeness Matters for RL]
The Banach fixed-point theorem (proved in the next lesson) states: if $T$ is a
contraction on a \emph{complete} metric space, then $T$ has a unique fixed point
and $T^n x \to x^*$ for any starting point $x$. The Bellman operator is a
$\gamma$-contraction on the Banach space of bounded functions with the sup norm.
Without completeness, we cannot guarantee that the sequence of value iterates
converges to anything in our space. This is the mathematical reason we work in
Banach spaces, not just normed spaces.
\end{rlconnection}

\subsection{Schauder Basis}
\label{subsec:schauder}

\begin{definition}[Schauder Basis]
\label{def:schauder}
A sequence $(e_n)_{n=1}^\infty$ in a Banach space $X$ is a \emph{Schauder basis}
if every $x \in X$ has a unique representation $x = \sum_{n=1}^\infty \alpha_n e_n$
(convergent in norm).
\end{definition}

\begin{example}
The standard unit vectors $(e_n)$ form a Schauder basis for $\ell^p$ ($1 \le p < \infty$):
every $x = (x_1, x_2, \ldots) \in \ell^p$ satisfies
$x = \sum_{n=1}^\infty x_n e_n$ in $\|\cdot\|_p$.
Note that $(e_n)$ is \emph{not} a Schauder basis for $\ell^\infty$, because the
partial sums $\sum_{n=1}^N x_n e_n$ need not converge to $x$ in $\|\cdot\|_\infty$.
\end{example}

%% ============================================================
%% SECTION 5: BOUNDED LINEAR OPERATORS
%% ============================================================
\section{Bounded Linear Operators}
\label{sec:operators}

\subsection{Definitions and Basic Properties}

\begin{definition}[Linear Operator]
\label{def:linear_op}
Let $X$ and $Y$ be normed spaces. A map $T \colon X \to Y$ is a \emph{linear
operator} if $T(\alpha x + \beta y) = \alpha T(x) + \beta T(y)$ for all
$x, y \in X$ and scalars $\alpha, \beta$.
\end{definition}

\begin{definition}[Bounded Linear Operator]
\label{def:bounded_op}
A linear operator $T \colon X \to Y$ is \emph{bounded} if there exists $c \ge 0$
such that $\|Tx\|_Y \le c \|x\|_X$ for all $x \in X$. The \emph{operator norm} is
\[
\|T\| = \sup_{x \ne 0} \frac{\|Tx\|_Y}{\|x\|_X} = \sup_{\|x\|_X = 1} \|Tx\|_Y
= \sup_{\|x\|_X \le 1} \|Tx\|_Y.
\]
\end{definition}

\begin{proposition}[Equivalent Formulations of the Operator Norm]
\label{prop:op_norm_equiv}
The three suprema in Definition~\ref{def:bounded_op} are equal.
\end{proposition}

\begin{proof}
Let $\alpha = \sup_{x \ne 0} \|Tx\|/\|x\|$, $\beta = \sup_{\|x\|=1} \|Tx\|$,
$\gamma = \sup_{\|x\| \le 1} \|Tx\|$.

$\alpha = \beta$: For $x \ne 0$, set $u = x/\|x\|$, so $\|u\| = 1$ and
$\|Tx\|/\|x\| = \|Tu\|$. Thus $\sup_{x \ne 0} \|Tx\|/\|x\| =
\sup_{\|u\|=1} \|Tu\|$.

$\beta \le \gamma$: The set $\{\|x\| = 1\}$ is contained in $\{\|x\| \le 1\}$.

$\gamma \le \alpha$: For $\|x\| \le 1$ with $x \ne 0$,
$\|Tx\| = \|x\| \cdot \|Tx\|/\|x\| \le 1 \cdot \alpha = \alpha$. Also $\|T(0)\| = 0 \le \alpha$.
So $\gamma \le \alpha = \beta \le \gamma$, giving equality.
\end{proof}

\begin{theorem}[Bounded $\Leftrightarrow$ Continuous]
\label{thm:bounded_iff_cts}
A linear operator $T \colon X \to Y$ between normed spaces is bounded if and only
if it is continuous.
\end{theorem}

\begin{proof}
($\Rightarrow$) If $T$ is bounded, then $\|Tx - Ty\| = \|T(x-y)\| \le \|T\| \|x - y\|$,
so $T$ is Lipschitz, hence continuous.

($\Leftarrow$) Suppose $T$ is continuous but not bounded. Then for each $n$, there
exists $x_n$ with $\|x_n\| = 1$ and $\|Tx_n\| > n$. Set $y_n = x_n/n$. Then
$\|y_n\| = 1/n \to 0$, so $y_n \to 0$, but $\|Ty_n\| = \|Tx_n\|/n > 1$. This
contradicts the continuity of $T$ at $0$ (since $y_n \to 0$ should imply
$Ty_n \to T(0) = 0$).
\end{proof}

\begin{remark}
For linear operators, continuity at a single point (say $0$) implies continuity
everywhere, because $T(x_n) - T(x) = T(x_n - x)$ and $x_n - x \to 0$.
\end{remark}

\begin{warning}[Unbounded Operators Exist]
Not every linear operator is bounded. On the space of polynomials $P[0,1]$ (with
the sup norm), the differentiation operator $D(f) = f'$ is linear but unbounded:
$\|x^n\|_\infty = 1$ but $\|D(x^n)\|_\infty = \|n x^{n-1}\|_\infty = n \to \infty$.
\end{warning}

\subsection{The Space $B(X,Y)$}

\begin{definition}[Space of Bounded Operators]
\label{def:BXY}
For normed spaces $X, Y$, we write $B(X,Y)$ for the set of all bounded linear
operators $T \colon X \to Y$, equipped with the operator norm. When $Y = X$, we
write $B(X) = B(X,X)$.
\end{definition}

\begin{theorem}[$B(X,Y)$ is a Normed Space]
\label{thm:BXY_normed}
$B(X,Y)$ is a normed space under the operator norm.
\end{theorem}

\begin{proof}
$B(X,Y)$ is a vector space: for $T, S \in B(X,Y)$ and scalar $\alpha$,
$(T + S)(x) = Tx + Sx$ and $(\alpha T)(x) = \alpha Tx$ are bounded linear operators
(with $\|T + S\| \le \|T\| + \|S\|$ and $\|\alpha T\| = |\alpha| \|T\|$).

The operator norm satisfies the norm axioms:
\begin{itemize}[nosep]
\item \ref{N1}: $\|T\| = 0 \iff \|Tx\| = 0$ for all $\|x\| \le 1 \iff Tx = 0$
for all $x \iff T = 0$.
\item \ref{N2}: $\|\alpha T\| = \sup_{\|x\| \le 1} \|\alpha Tx\| =
|\alpha| \sup_{\|x\| \le 1} \|Tx\| = |\alpha| \|T\|$.
\item \ref{N3}: $\|T + S\| = \sup_{\|x\| \le 1} \|(T+S)x\| \le
\sup_{\|x\| \le 1} (\|Tx\| + \|Sx\|) \le \|T\| + \|S\|$.
\end{itemize}
\end{proof}

\begin{theorem}[$B(X,Y)$ is Banach When $Y$ is Banach]
\label{thm:BXY_banach}
If $Y$ is a Banach space, then $B(X,Y)$ is a Banach space.
\end{theorem}

\begin{proof}
Let $(T_n)$ be a Cauchy sequence in $B(X,Y)$. For each $x \in X$,
$\|T_n x - T_m x\| \le \|T_n - T_m\| \|x\| \to 0$, so $(T_n x)$ is Cauchy in $Y$.
Since $Y$ is Banach, the limit exists: define $Tx = \lim_{n \to \infty} T_n x$.

\textbf{$T$ is linear.} $T(\alpha x + \beta y) = \lim T_n(\alpha x + \beta y) =
\lim (\alpha T_n x + \beta T_n y) = \alpha Tx + \beta Ty$.

\textbf{$T$ is bounded.} Since $(T_n)$ is Cauchy, it is bounded: $\|T_n\| \le M$
for some $M$. For each $x$, $\|Tx\| = \lim \|T_n x\| \le M \|x\|$, so
$\|T\| \le M$.

\textbf{$T_n \to T$ in operator norm.} Given $\varepsilon > 0$, choose $N$ such
that $\|T_n - T_m\| < \varepsilon$ for $n, m \ge N$. For any $x$ with $\|x\| \le 1$
and $n \ge N$:
$\|T_n x - T_m x\| \le \varepsilon$ for all $m \ge N$. Letting $m \to \infty$:
$\|T_n x - Tx\| \le \varepsilon$. Taking the sup over $\|x\| \le 1$:
$\|T_n - T\| \le \varepsilon$ for $n \ge N$.
\end{proof}

\subsection{Important Examples}

\begin{example}[Matrix Operators]
\label{ex:matrix_op}
Every $m \times n$ real matrix $A$ defines a bounded linear operator
$T_A \colon (\R^n, \|\cdot\|_p) \to (\R^m, \|\cdot\|_p)$ by $T_A(x) = Ax$.
Boundedness follows from finite-dimensionality (all linear operators on
finite-dimensional spaces are bounded, as we prove below).
\end{example}

\begin{proposition}[All Linear Operators on Finite-Dimensional Spaces are Bounded]
\label{prop:finite_dim_bounded}
If $X$ is a finite-dimensional normed space and $Y$ is any normed space, then
every linear operator $T \colon X \to Y$ is bounded.
\end{proposition}

\begin{proof}
Let $\{e_1, \ldots, e_n\}$ be a basis for $X$. For $x = \sum_{i=1}^n \alpha_i e_i$,
by the equivalence of norms (Theorem~\ref{thm:norm_equiv}), there exists $c > 0$
with $\sum |\alpha_i| \le c \|x\|_X$ (since $x \mapsto \sum |\alpha_i|$ is a norm
on $X$, equivalent to $\|\cdot\|_X$). Then:
\[
\|Tx\| = \Bigl\|\sum_{i=1}^n \alpha_i Te_i\Bigr\| \le \sum_{i=1}^n |\alpha_i| \|Te_i\|
\le c \|x\|_X \cdot \max_i \|Te_i\|.
\]
So $T$ is bounded with $\|T\| \le c \cdot \max_i \|Te_i\|$.
\end{proof}

\begin{example}[Integral Operator]
\label{ex:integral_op}
Let $k \in C([a,b] \times [a,b])$ and define $T \colon C[a,b] \to C[a,b]$ by
\[
(Tf)(s) = \int_a^b k(s,t)\, f(t)\, dt.
\]
This is a bounded linear operator with
$\|T\| \le \max_{s \in [a,b]} \int_a^b |k(s,t)|\, dt$.
Indeed, $|(Tf)(s)| \le \int_a^b |k(s,t)| |f(t)|\, dt \le \|f\|_\infty \int_a^b |k(s,t)|\, dt$.
\end{example}

\begin{example}[Right Shift Operator]
\label{ex:shift}
On $\ell^p$ ($1 \le p \le \infty$), the right shift $R(x_1, x_2, \ldots) =
(0, x_1, x_2, \ldots)$ is a bounded linear operator with $\|R\| = 1$: clearly
$\|Rx\|_p = \|x\|_p$ for all $x$.
\end{example}

\begin{rlconnection}[The Bellman Operator is Bounded]
The Bellman expectation operator $T^\pi \colon b(\mathcal{S}) \to b(\mathcal{S})$
defined by $(T^\pi V)(s) = r^\pi(s) + \gamma \sum_{s'} P^\pi(s'|s) V(s')$ is a
bounded (in fact, affine) operator on the Banach space of bounded functions. Its
``linear part'' $V \mapsto \gamma \sum_{s'} P^\pi(s'|s) V(s')$ has operator norm
at most $\gamma < 1$ (under $\|\cdot\|_\infty$), which is why $T^\pi$ is a
contraction. The framework of $B(X,Y)$ and operator norms gives us the precise
language to state and prove this.
\end{rlconnection}

\subsection{Dual Spaces}

\begin{definition}[Dual Space]
\label{def:dual}
The \emph{(continuous) dual space} of a normed space $X$ is
$X^* = B(X, \R)$ (or $B(X, \C)$ if $X$ is a complex space): the space of all
bounded linear functionals on $X$, equipped with the operator norm
$\|f\|_{X^*} = \sup_{\|x\| \le 1} |f(x)|$.
\end{definition}

\begin{corollary}
\label{cor:dual_banach}
For any normed space $X$, the dual space $X^*$ is a Banach space.
\end{corollary}

\begin{proof}
$\R$ (or $\C$) is a Banach space, so $X^* = B(X, \R)$ is Banach by
Theorem~\ref{thm:BXY_banach}.
\end{proof}

\begin{remark}
The dual space is always complete, even if $X$ itself is not.
\end{remark}

\begin{example}[Dual of $\ell^p$]
\label{ex:dual_lp}
For $1 \le p < \infty$ with $1/p + 1/q = 1$ (where $q = \infty$ when $p = 1$),
$(\ell^p)^* \cong \ell^q$ isometrically: every bounded linear functional on
$\ell^p$ has the form $f(x) = \sum_{n=1}^\infty y_n x_n$ for a unique
$y = (y_n) \in \ell^q$, with $\|f\| = \|y\|_q$. This is the Riesz representation
theorem for $\ell^p$ spaces.
\end{example}

%% ============================================================
%% SECTION 6: COMPACTNESS
%% ============================================================
\section{Compactness}
\label{sec:compactness}

Compactness is the key tool for passing from local to global properties and for
establishing existence results. In infinite-dimensional spaces, compactness
behaves very differently from the finite-dimensional case.

\subsection{Definitions and Basic Theory}

\begin{definition}[Compact Set]
\label{def:compact}
A subset $K$ of a metric space $(X,d)$ is \emph{compact} if every open cover of
$K$ has a finite subcover: whenever $K \subseteq \bigcup_{\alpha \in I} G_\alpha$
with each $G_\alpha$ open, there exist finitely many $\alpha_1, \ldots, \alpha_n$
with $K \subseteq G_{\alpha_1} \cup \cdots \cup G_{\alpha_n}$.
\end{definition}

\begin{definition}[Sequential Compactness]
\label{def:seq_compact}
A subset $K$ of a metric space is \emph{sequentially compact} if every sequence in
$K$ has a subsequence that converges to a point in $K$.
\end{definition}

\begin{definition}[Totally Bounded]
\label{def:totally_bounded}
A subset $A$ of a metric space $(X,d)$ is \emph{totally bounded} if for every
$\varepsilon > 0$, $A$ can be covered by finitely many open balls of radius
$\varepsilon$: there exist $x_1, \ldots, x_n \in X$ with
$A \subseteq \bigcup_{i=1}^n B(x_i, \varepsilon)$.
\end{definition}

\begin{proposition}[Totally Bounded Implies Bounded]
\label{prop:tb_bounded}
Every totally bounded set is bounded. The converse fails in general.
\end{proposition}

\begin{proof}
If $A \subseteq \bigcup_{i=1}^n B(x_i, 1)$, then for any $a, b \in A$, $a \in B(x_i, 1)$
and $b \in B(x_j, 1)$ for some $i, j$, so $d(a,b) \le d(a, x_i) + d(x_i, x_j) +
d(x_j, b) < 1 + \max_{i,j} d(x_i, x_j) + 1$. Hence $\diam(A) < \infty$.

For the failure of the converse: the closed unit ball in $\ell^2$ is bounded but not
totally bounded (the standard basis vectors $e_n$ satisfy $\|e_n - e_m\| = \sqrt{2}$
for $n \ne m$, so no finite collection of balls of radius $\sqrt{2}/2$ can cover them all).
\end{proof}

\begin{theorem}[Characterization of Compactness in Metric Spaces]
\label{thm:compact_equiv}
For a subset $K$ of a metric space $(X,d)$, the following are equivalent:
\begin{enumerate}[nosep,label=(\roman*)]
    \item $K$ is compact.
    \item $K$ is sequentially compact.
    \item $K$ is complete and totally bounded.
\end{enumerate}
\end{theorem}

\begin{proof}
\textbf{(i) $\Rightarrow$ (ii):}
Let $(x_n)$ be a sequence in $K$. Suppose for contradiction that $(x_n)$ has no
convergent subsequence in $K$. Then every point $y \in K$ has a neighborhood
$B(y, r_y)$ containing only finitely many terms of $(x_n)$ (otherwise, we could
extract a subsequence converging to $y$). The collection $\{B(y, r_y) : y \in K\}$
covers $K$, so by compactness there is a finite subcover. But a finite union of
sets, each containing finitely many terms, contains only finitely many terms total,
contradicting the fact that $(x_n)$ has infinitely many terms in $K$.

\textbf{(ii) $\Rightarrow$ (iii):}
\emph{Completeness:} Let $(x_n)$ be Cauchy in $K$. By sequential compactness,
$(x_n)$ has a subsequence converging to some $x \in K$. A Cauchy sequence with a
convergent subsequence converges (Theorem~\ref{thm:R_complete} proof argument), so
$x_n \to x \in K$.

\emph{Total boundedness:} Suppose $K$ is not totally bounded. Then there exists
$\varepsilon > 0$ such that $K$ cannot be covered by finitely many
$\varepsilon$-balls. Pick $x_1 \in K$ arbitrarily. Having chosen $x_1, \ldots, x_n$,
since $K \not\subseteq \bigcup_{i=1}^n B(x_i, \varepsilon)$, choose
$x_{n+1} \in K \setminus \bigcup_{i=1}^n B(x_i, \varepsilon)$. Then
$d(x_{n+1}, x_i) \ge \varepsilon$ for all $i \le n$. The sequence $(x_n)$ has no
convergent subsequence (any two terms are at least $\varepsilon$ apart),
contradicting sequential compactness.

\textbf{(iii) $\Rightarrow$ (i):}
Suppose $K$ is complete and totally bounded. We use the \emph{Lebesgue covering
lemma} approach.

\emph{Claim (Lebesgue number):} For every open cover $\{G_\alpha\}$ of $K$, there
exists $\delta > 0$ (the \emph{Lebesgue number}) such that for every $x \in K$,
$B(x, \delta) \subseteq G_\alpha$ for some $\alpha$.

\emph{Proof of claim:} Suppose not. Then for each $n$, there exists $x_n \in K$
with $B(x_n, 1/n) \not\subseteq G_\alpha$ for any $\alpha$. By (iii) $\Rightarrow$
(ii) (which we now verify: complete + totally bounded $\Rightarrow$ sequentially
compact, see below), $(x_n)$ has a subsequence $x_{n_k} \to x \in K$. Since
$x \in K \subseteq \bigcup_\alpha G_\alpha$, we have $x \in G_\beta$ for some
$\beta$. Since $G_\beta$ is open, $B(x, r) \subseteq G_\beta$ for some $r > 0$.
For large $k$, $d(x_{n_k}, x) < r/2$ and $1/n_k < r/2$, so
$B(x_{n_k}, 1/n_k) \subseteq B(x, r) \subseteq G_\beta$, contradicting the choice
of $x_{n_k}$.

To complete the argument, we need: complete + totally bounded implies sequentially
compact. Let $(x_n)$ be a sequence in $K$. Cover $K$ by finitely many balls of
radius $1$; one ball contains infinitely many terms, giving a subsequence
$(x_n^{(1)})$ with $\diam \le 2$. Cover $K$ by balls of radius $1/2$; one contains
infinitely many terms of $(x_n^{(1)})$, giving $(x_n^{(2)})$ with $\diam \le 1$.
Continue: at stage $k$, get $(x_n^{(k)})$ contained in a ball of radius $1/k$.
The diagonal subsequence $y_k = x_k^{(k)}$ is Cauchy (for $k, m \ge N$,
$d(y_k, y_m) \le 2/N$), hence converges by completeness.

Now: given the Lebesgue number $\delta > 0$, cover $K$ by finitely many balls
of radius $\delta/2$ (by total boundedness). Each ball has diameter $\le \delta$,
so is contained in some $G_\alpha$. Taking the corresponding finitely many
$G_\alpha$'s gives a finite subcover.
\end{proof}

\begin{corollary}[Heine--Borel for $\R^n$]
\label{cor:heine_borel}
A subset of $\R^n$ is compact if and only if it is closed and bounded.
\end{corollary}

\begin{proof}
($\Rightarrow$) Compact subsets of metric spaces are closed and bounded.
(\emph{Closed:} If $x_n \in K$ and $x_n \to x$, then by sequential compactness
there is a subsequence converging in $K$; by uniqueness of limits, $x \in K$.
\emph{Bounded:} Compact implies totally bounded implies bounded.)

($\Leftarrow$) Let $K \subseteq \R^n$ be closed and bounded. Bounded in $\R^n$
implies totally bounded (cover $K$ by a grid of cubes of side $\varepsilon/\sqrt{n}$).
Closed subset of a complete space is complete. By Theorem~\ref{thm:compact_equiv},
$K$ is compact.
\end{proof}

\subsection{Riesz's Lemma and Compactness in Infinite Dimensions}

\begin{lemma}[Riesz's Lemma]
\label{lem:riesz}
Let $Y$ be a proper closed subspace of a normed space $X$ (i.e., $Y \ne X$). For
every $\theta \in (0,1)$, there exists $x_\theta \in X$ with $\|x_\theta\| = 1$
and $d(x_\theta, Y) = \inf_{y \in Y} \|x_\theta - y\| \ge \theta$.
\end{lemma}

\begin{proof}
Since $Y \ne X$, there exists $z \in X \setminus Y$. Since $Y$ is closed,
$d = d(z, Y) = \inf_{y \in Y} \|z - y\| > 0$. Choose $y_0 \in Y$ with
$\|z - y_0\| < d/\theta$ (possible since $d/\theta > d = \inf$). Define
\[
x_\theta = \frac{z - y_0}{\|z - y_0\|}.
\]
Then $\|x_\theta\| = 1$. For any $y \in Y$:
\[
\|x_\theta - y\| = \frac{\|z - y_0 - \|z - y_0\| y\|}{\|z - y_0\|}
= \frac{\|z - \underbrace{(y_0 + \|z - y_0\| y)}_{\in Y}\|}{\|z - y_0\|}
\ge \frac{d}{\|z - y_0\|} > \frac{d}{d/\theta} = \theta.
\]
Taking the infimum over $y \in Y$: $d(x_\theta, Y) \ge \theta$.
\end{proof}

\begin{remark}
We cannot always achieve $\theta = 1$ (i.e., exact distance $1$ from $Y$). This is
possible in reflexive spaces but not in general.
\end{remark}

\begin{theorem}[Riesz's Theorem: Compactness Characterizes Finite Dimension]
\label{thm:riesz_compact}
The closed unit ball $\overline{B}(0,1) = \{x \in X : \|x\| \le 1\}$ in a normed
space $X$ is compact if and only if $X$ is finite-dimensional.
\end{theorem}

\begin{proof}
($\Leftarrow$) If $\dim X = n < \infty$, then $X$ is isomorphic to $\R^n$
(Theorem~\ref{thm:norm_equiv}), and the closed unit ball is closed and bounded,
hence compact by Heine--Borel (Corollary~\ref{cor:heine_borel}).

($\Rightarrow$) We prove the contrapositive: if $X$ is infinite-dimensional, then
$\overline{B}(0,1)$ is not compact (not sequentially compact). We construct a
sequence in $\overline{B}(0,1)$ with no convergent subsequence.

Pick $x_1 \in X$ with $\|x_1\| = 1$. Let $Y_1 = \Span\{x_1\}$, which is a proper
closed subspace (finite-dimensional subspaces are closed). By Riesz's lemma with
$\theta = 1/2$, there exists $x_2$ with $\|x_2\| = 1$ and $\|x_2 - y\| \ge 1/2$
for all $y \in Y_1$; in particular $\|x_2 - x_1\| \ge 1/2$.

Let $Y_2 = \Span\{x_1, x_2\}$, another proper closed subspace. By Riesz's lemma,
there exists $x_3$ with $\|x_3\| = 1$ and $d(x_3, Y_2) \ge 1/2$. Continuing
inductively, we obtain a sequence $(x_n)$ in $\overline{B}(0,1)$ with
$\|x_n - x_m\| \ge 1/2$ for all $n \ne m$.

No subsequence of $(x_n)$ can be Cauchy (hence none can converge), so
$\overline{B}(0,1)$ is not sequentially compact, hence not compact.
\end{proof}

\begin{keyresult}[Riesz's Theorem]
The closed unit ball in a normed space is compact if and only if the space is
finite-dimensional. This is one of the most important structural results in
functional analysis: it says that in infinite-dimensional spaces, bounded sets are
never compact (in the norm topology). This forces us to use weaker topologies or
stronger conditions (like equicontinuity) to obtain compactness.
\end{keyresult}

\begin{rlconnection}[Compactness and Function Approximation]
Riesz's theorem explains a fundamental challenge in RL: the space of all value
functions on an infinite state space is infinite-dimensional, so the unit ball is
not compact. This means that a bounded sequence of value function approximations
need not have a convergent subsequence. In practice, this is why we restrict to
finite-dimensional parameterizations (linear function approximation, neural
networks) or use compactness criteria like Arzel\`{a}--Ascoli
(Theorem~\ref{thm:arzela_ascoli}) to guarantee that approximation sequences converge.
\end{rlconnection}

\subsection{The Arzel\`{a}--Ascoli Theorem}

\begin{definition}[Equicontinuity]
\label{def:equicontinuous}
A family $\mathcal{F} \subseteq C[a,b]$ is \emph{equicontinuous} if for every
$\varepsilon > 0$ there exists $\delta > 0$ such that
$|f(s) - f(t)| < \varepsilon$ for all $f \in \mathcal{F}$ and all $s, t \in [a,b]$
with $|s - t| < \delta$. The family is \emph{uniformly bounded} if
$\sup_{f \in \mathcal{F}} \|f\|_\infty < \infty$.
\end{definition}

\begin{theorem}[Arzel\`{a}--Ascoli]
\label{thm:arzela_ascoli}
A subset $\mathcal{F} \subseteq C[a,b]$ (with the sup norm) is relatively compact
(i.e., $\cl(\mathcal{F})$ is compact) if and only if $\mathcal{F}$ is uniformly
bounded and equicontinuous.
\end{theorem}

\begin{proof}
($\Leftarrow$) \textbf{Uniformly bounded + equicontinuous implies relatively compact.}
By Theorem~\ref{thm:compact_equiv}, it suffices to show that every sequence $(f_n)$
in $\mathcal{F}$ has a uniformly convergent subsequence (the closure inherits
completeness from $C[a,b]$).

\textbf{Step 1: Pointwise convergence on a dense set.}
Let $D = \{r_1, r_2, \ldots\}$ be a countable dense subset of $[a,b]$ (e.g., the
rationals in $[a,b]$). The sequence $(f_n(r_1))$ is bounded (by uniform
boundedness), so by Bolzano--Weierstrass it has a convergent subsequence
$(f_{n_k^{(1)}}(r_1))$. From this subsequence, extract a further subsequence
convergent at $r_2$: $(f_{n_k^{(2)}})$. Continue for each $r_j$. The diagonal
subsequence $g_k = f_{n_k^{(k)}}$ converges at every $r_j \in D$.

\textbf{Step 2: Uniform convergence.}
Given $\varepsilon > 0$, choose $\delta > 0$ from equicontinuity for
$\varepsilon/3$. Cover $[a,b]$ by finitely many intervals of length $< \delta$, and
pick a point $r_{j_i} \in D$ in each interval, giving $r_{j_1}, \ldots, r_{j_m}$.
Since $(g_k)$ converges at each $r_{j_i}$, choose $K$ such that
$|g_k(r_{j_i}) - g_l(r_{j_i})| < \varepsilon/3$ for all $k, l \ge K$ and all $i$.

For any $t \in [a,b]$, choose $r_{j_i}$ with $|t - r_{j_i}| < \delta$. Then for
$k, l \ge K$:
\begin{align*}
|g_k(t) - g_l(t)| &\le |g_k(t) - g_k(r_{j_i})| + |g_k(r_{j_i}) - g_l(r_{j_i})|
+ |g_l(r_{j_i}) - g_l(t)| \\
&< \varepsilon/3 + \varepsilon/3 + \varepsilon/3 = \varepsilon.
\end{align*}
So $\|g_k - g_l\|_\infty < \varepsilon$ for $k, l \ge K$, showing $(g_k)$ is Cauchy
in $(C[a,b], \|\cdot\|_\infty)$, hence convergent.

($\Rightarrow$) \textbf{Relatively compact implies uniformly bounded and equicontinuous.}

\emph{Uniformly bounded:} If $\cl(\mathcal{F})$ is compact, it is bounded, so
$\sup_{f \in \mathcal{F}} \|f\|_\infty \le \sup_{g \in \cl(\mathcal{F})} \|g\|_\infty
< \infty$.

\emph{Equicontinuous:} Given $\varepsilon > 0$, cover $\cl(\mathcal{F})$ by finitely
many balls $B(g_i, \varepsilon/3)$, $i = 1, \ldots, N$. Each $g_i \in C[a,b]$ is
uniformly continuous on $[a,b]$, so choose $\delta_i > 0$ with
$|g_i(s) - g_i(t)| < \varepsilon/3$ when $|s-t| < \delta_i$. Set
$\delta = \min(\delta_1, \ldots, \delta_N) > 0$. For any $f \in \mathcal{F}$,
choose $g_i$ with $\|f - g_i\|_\infty < \varepsilon/3$. For $|s-t| < \delta$:
$|f(s) - f(t)| \le |f(s) - g_i(s)| + |g_i(s) - g_i(t)| + |g_i(t) - f(t)| <
\varepsilon/3 + \varepsilon/3 + \varepsilon/3 = \varepsilon$.
\end{proof}

\begin{rlconnection}[Arzel\`{a}--Ascoli in Policy Optimization]
In continuous-state RL, when we optimize over a class of value functions or
policies, we need compactness to guarantee that maximizing sequences have
convergent subsequences. The Arzel\`{a}--Ascoli theorem provides the criterion:
if value function approximations are uniformly bounded and equicontinuous (e.g.,
Lipschitz with a common Lipschitz constant), then any sequence has a uniformly
convergent subsequence. This is used in the analysis of fitted value iteration
and actor-critic methods on continuous state spaces.
\end{rlconnection}

\subsection{Compact Operators}

\begin{definition}[Compact Operator]
\label{def:compact_op}
A linear operator $T \colon X \to Y$ between normed spaces is \emph{compact} if
$T$ maps bounded sets to relatively compact sets, i.e., for every bounded
$B \subseteq X$, the closure $\cl(T(B))$ is compact in $Y$.
Equivalently, $T$ is compact if every bounded sequence $(x_n)$ in $X$ has a
subsequence $(x_{n_k})$ such that $(Tx_{n_k})$ converges in $Y$.
\end{definition}

\begin{proposition}[Compact Operators are Bounded]
\label{prop:compact_bounded}
Every compact operator is bounded.
\end{proposition}

\begin{proof}
Suppose $T$ is compact but not bounded. Then there exist $(x_n)$ with $\|x_n\| = 1$
and $\|Tx_n\| \to \infty$. Since $(x_n)$ is bounded, compactness gives a
subsequence with $(Tx_{n_k})$ convergent, hence bounded. But $\|Tx_{n_k}\| \to
\infty$, a contradiction.
\end{proof}

\begin{example}[Integral Operators are Often Compact]
\label{ex:compact_integral}
The integral operator from Example~\ref{ex:integral_op} with continuous kernel
$k$ on $[a,b] \times [a,b]$ is compact as an operator $C[a,b] \to C[a,b]$. This
follows from Arzel\`{a}--Ascoli: if $(f_n)$ is bounded in $C[a,b]$, then
$(Tf_n)$ is uniformly bounded (since $T$ is bounded) and equicontinuous (since $k$
is uniformly continuous on the compact set $[a,b]^2$, and
$|(Tf_n)(s_1) - (Tf_n)(s_2)| \le \|f_n\|_\infty \int_a^b |k(s_1,t) - k(s_2,t)|\, dt$
is small uniformly in $n$ when $|s_1 - s_2|$ is small).
\end{example}

%% ============================================================
%% SECTION 7: BAIRE CATEGORY THEOREM
%% ============================================================
\section{The Baire Category Theorem}
\label{sec:baire}

The Baire category theorem is one of the most powerful existence results in
analysis. It underpins the three fundamental theorems of Banach space theory:
the Uniform Boundedness Principle, the Open Mapping Theorem, and the Closed Graph
Theorem (all treated in the next lesson).

\subsection{Nowhere Dense Sets and Category}

\begin{definition}[Nowhere Dense]
\label{def:nowhere_dense}
A subset $A$ of a metric space $X$ is \emph{nowhere dense} if $\Int(\cl(A)) =
\emptyset$, i.e., the closure of $A$ has empty interior. Equivalently, $A$ is
nowhere dense if and only if for every non-empty open set $U$, there is a
non-empty open set $V \subseteq U$ with $V \cap A = \emptyset$.
\end{definition}

\begin{example}
In $\R$: $\Z$ is nowhere dense, $\{1/n : n \in \N\}$ is nowhere dense, $\Q$ is
\emph{not} nowhere dense (its closure is $\R$, which has non-empty interior).
The Cantor set is a closed, uncountable, nowhere dense subset of $[0,1]$.
\end{example}

\begin{definition}[First and Second Category]
\label{def:category}
A subset of a metric space is of \emph{first category} (or \emph{meager}) if it is
a countable union of nowhere dense sets. A set is of \emph{second category} (or
\emph{non-meager}) if it is not of first category.
\end{definition}

\begin{remark}
``First category'' means ``small'' in the topological sense, and ``second category''
means ``large.'' This is a different notion of smallness from measure zero.
\end{remark}

\subsection{Statement and Proof}

\begin{theorem}[Baire Category Theorem]
\label{thm:baire}
Every complete metric space is of second category in itself. Equivalently, if
$(X,d)$ is a complete metric space and $X = \bigcup_{n=1}^\infty F_n$ with each
$F_n$ closed, then at least one $F_n$ has non-empty interior.
\end{theorem}

\begin{proof}
We prove the equivalent contrapositive: if $A_1, A_2, \ldots$ are nowhere dense
subsets of a complete metric space $X$, then $\bigcup_{n=1}^\infty A_n \ne X$.

Since $A_1$ is nowhere dense, $X \setminus \cl(A_1)$ is a non-empty open set (if
$\cl(A_1) = X$, then $\Int(\cl(A_1)) = X \ne \emptyset$, contradicting nowhere
density). Choose $x_1 \in X \setminus \cl(A_1)$ and $r_1 > 0$ with
$\overline{B}(x_1, r_1) \subseteq X \setminus \cl(A_1)$ and $r_1 < 1$. In
particular, $\overline{B}(x_1, r_1) \cap A_1 = \emptyset$.

Since $A_2$ is nowhere dense and $B(x_1, r_1)$ is a non-empty open set, there
exists a non-empty open ball $B(x_2, r_2) \subseteq B(x_1, r_1)$ with
$\overline{B}(x_2, r_2) \cap \cl(A_2) = \emptyset$ and $r_2 < 1/2$.

Continuing inductively, at step $n$ we find $x_n$ and $0 < r_n < 1/n$ with
\[
\overline{B}(x_n, r_n) \subseteq B(x_{n-1}, r_{n-1}) \quad \text{and} \quad
\overline{B}(x_n, r_n) \cap A_n = \emptyset.
\]

The sequence $(x_n)$ is Cauchy: for $m > n$, $x_m \in B(x_n, r_n)$, so
$d(x_m, x_n) < r_n < 1/n \to 0$. Since $X$ is complete, $x_n \to x$ for some
$x \in X$.

For each $n$, $x_m \in \overline{B}(x_n, r_n)$ for all $m \ge n$ (by nesting), so
$x = \lim x_m \in \overline{B}(x_n, r_n)$ (closed balls are closed). Therefore
$x \notin A_n$ for any $n$, so $x \notin \bigcup_{n=1}^\infty A_n$. Hence
$\bigcup_{n=1}^\infty A_n \ne X$.
\end{proof}

\begin{corollary}
\label{cor:baire_dense}
In a complete metric space, the intersection of countably many dense open sets is
dense (in fact, it is a dense $G_\delta$ set).
\end{corollary}

\begin{proof}
Let $G_1, G_2, \ldots$ be dense open sets in a complete metric space $X$. Their
complements $F_n = X \setminus G_n$ are closed and nowhere dense (since $G_n$ is
dense, $\cl(F_n) = F_n$ has empty interior: any ball $B(x,r)$ contains a point of
$G_n$, hence a point outside $F_n$).

By the Baire category theorem, $\bigcup F_n \ne X$, so
$\bigcap G_n = X \setminus \bigcup F_n \ne \emptyset$.

To show density: let $U$ be any non-empty open set. Then $U$ is a complete metric
space (as an open subset of $X$... actually we need a slight modification). Instead,
apply the construction in the proof of Theorem~\ref{thm:baire} starting from any
non-empty open ball $B(x_0, r_0) \subseteq U$: since each $G_n$ is dense,
$B(x_0, r_0) \cap G_n$ is a non-empty open set. Choose nested closed balls
$\overline{B}(x_n, r_n) \subseteq B(x_{n-1}, r_{n-1}) \cap G_n$ with $r_n \to 0$.
The limit $x = \lim x_n$ lies in $\overline{B}(x_0, r_0) \subseteq U$ and in each
$G_n$, so $x \in U \cap \bigcap G_n$.
\end{proof}

\begin{corollary}[$\Q$ is Not a $G_\delta$ in $\R$]
\label{cor:Q_not_Gdelta}
The rationals $\Q$ cannot be written as a countable intersection of open sets.
\end{corollary}

\begin{proof}
Suppose $\Q = \bigcap_{n=1}^\infty G_n$ with each $G_n$ open. Since $\Q$ is dense,
each $G_n$ is dense. Also, for each $q \in \Q$, $\R \setminus \{q\}$ is open and
dense. By Corollary~\ref{cor:baire_dense}, $\bigcap_{n} G_n \cap
\bigcap_{q \in \Q} (\R \setminus \{q\}) = \Q \cap (\R \setminus \Q) = \emptyset$
is dense, which is a contradiction (the empty set is not dense in $\R$).
\end{proof}

\subsection{Application Preview}

\begin{remark}[Preview: The Uniform Boundedness Principle]
\label{rem:ubp_preview}
The Baire category theorem has the following powerful consequence (proved in the
next lesson): if $\{T_\alpha\}_{\alpha \in I}$ is a family of bounded linear
operators from a Banach space $X$ to a normed space $Y$ such that
$\sup_\alpha \|T_\alpha x\| < \infty$ for every $x \in X$ (pointwise boundedness),
then $\sup_\alpha \|T_\alpha\| < \infty$ (uniform boundedness). The proof works by
writing $X = \bigcup_{n=1}^\infty \{x : \sup_\alpha \|T_\alpha x\| \le n\}$ and
applying Baire to conclude that one of these closed sets has non-empty interior.
\end{remark}

\begin{rlconnection}[Baire Category and Genericity in RL]
The Baire category theorem is used in RL theory to prove ``genericity'' results:
that certain properties hold for ``most'' value functions or policies in a
topological sense. For instance, one can show that for a generic MDP, the optimal
policy is unique (the set of MDPs with multiple optimal policies is meager). The
Uniform Boundedness Principle, which follows from Baire, is used in the analysis
of temporal-difference learning to conclude that pointwise convergence of TD
estimates implies uniform bounds on approximation error.
\end{rlconnection}

%% ============================================================
%% SECTION 8: EXERCISES
%% ============================================================
\section{Exercises}
\label{sec:exercises}

\begin{exercise}[Metrics from Bounded Metrics]
\label{ex:bounded_metric}
Let $(X,d)$ be a metric space. Show that $\tilde{d}(x,y) = \min(d(x,y), 1)$ and
$\hat{d}(x,y) = d(x,y)/(1 + d(x,y))$ are both metrics on $X$ that generate the
same topology as $d$. Show that both $\tilde{d}$ and $\hat{d}$ are bounded metrics,
and that $(X,d)$ is complete if and only if $(X,\hat{d})$ is complete.
\end{exercise}

\begin{exercise}[Non-Completeness of $C[0,1]$ Under $L^1$]
\label{ex:C01_L1}
Give a detailed proof that $(C[0,1], \|\cdot\|_1)$ is not complete by
constructing an explicit Cauchy sequence that does not converge. (You may modify
the construction in Example~\ref{ex:Cab_not_banach}, but provide full details.)
\end{exercise}

\begin{exercise}[Closed Subspace of a Banach Space]
\label{ex:closed_subspace}
Let $X$ be a Banach space and $Y \subseteq X$ a subspace (i.e., a vector subspace).
Show that $Y$ is a Banach space (under the restricted norm) if and only if $Y$ is
a closed subset of $X$.
\end{exercise}

\begin{exercise}[Operator Norm Computation]
\label{ex:op_norm}
Let $T \colon (\R^2, \|\cdot\|_1) \to (\R^2, \|\cdot\|_\infty)$ be defined by
$T(x_1, x_2) = (x_1 + x_2, x_1 - x_2)$. Compute $\|T\|$ exactly.

\noindent\textit{Hint:} Use the definition $\|T\| = \sup_{\|x\|_1 \le 1} \|Tx\|_\infty$
and parametrize the unit ball in $\|\cdot\|_1$.
\end{exercise}

\begin{exercise}[Completeness of $\ell^\infty$]
\label{ex:ell_infty}
Prove directly (without appealing to Theorem~\ref{thm:ell_p_banach}) that
$(\ell^\infty, \|\cdot\|_\infty)$ is a Banach space.
\end{exercise}

\begin{exercise}[The Bellman Operator Norm]
\label{ex:bellman_norm}
Let $\mathcal{S} = \{1, \ldots, n\}$ be a finite state space and
$P \in \R^{n \times n}$ a stochastic matrix (non-negative entries, rows sum to $1$).
Define the linear operator $L \colon (\R^n, \|\cdot\|_\infty) \to
(\R^n, \|\cdot\|_\infty)$ by $(LV)(s) = \gamma \sum_{s'} P(s'|s) V(s')$ for
$\gamma \in (0,1)$.
\begin{enumerate}[nosep,label=(\alph*)]
    \item Show that $\|L\| = \gamma$.
    \item Conclude that $T = r + L$ (the Bellman operator for a fixed policy) is a
    contraction with factor $\gamma$ under $\|\cdot\|_\infty$.
    \item Verify the contraction property directly by computing
    $\|TV - TW\|_\infty$ for arbitrary $V, W \in \R^n$.
\end{enumerate}
\end{exercise}

\begin{exercise}[Riesz's Lemma is Sharp]
\label{ex:riesz_sharp}
Show that in $\ell^2$, Riesz's lemma can be strengthened to $\theta = 1$: for any
proper closed subspace $Y$, there exists $x$ with $\|x\| = 1$ and $d(x, Y) = 1$.

\noindent\textit{Hint:} Use the orthogonal projection onto $Y$.
\end{exercise}

\begin{exercise}[Equicontinuity and Compactness]
\label{ex:equicont}
Let $K > 0$ and $\mathcal{F}_K = \{f \in C[0,1] : |f(s) - f(t)| \le K|s-t|
\text{ for all } s, t, \text{ and } |f(0)| \le 1\}$.
\begin{enumerate}[nosep,label=(\alph*)]
    \item Show that $\mathcal{F}_K$ is equicontinuous and uniformly bounded.
    \item Conclude by Arzel\`{a}--Ascoli that $\mathcal{F}_K$ is relatively compact
    in $(C[0,1], \|\cdot\|_\infty)$.
    \item Show that $\mathcal{F}_K$ is in fact closed, hence compact.
\end{enumerate}
\end{exercise}

\begin{exercise}[Baire Category Application]
\label{ex:baire_app}
Use the Baire category theorem to show that $\R$ cannot be written as a countable
union of sets, each of which is closed and has empty interior.

\noindent Give an example showing that $\R$ \emph{can} be written as a countable
union of closed sets (some of which necessarily have non-empty interior).
\end{exercise}

\begin{exercise}[Nowhere Differentiable Functions are Generic]
\label{ex:nowhere_diff}
This exercise outlines the classical Baire category proof that ``most'' continuous
functions are nowhere differentiable.

Let $A_n = \{f \in C[0,1] : \text{there exists } x \in [0,1] \text{ with }
|f(x+h) - f(x)| \le n|h| \text{ for all sufficiently small } |h|\}$.
\begin{enumerate}[nosep,label=(\alph*)]
    \item Show that $A_n$ is a closed subset of $(C[0,1], \|\cdot\|_\infty)$.
    \item Show that $A_n$ has empty interior.
    \textit{Hint:} Approximate any $f \in A_n$ by a piecewise-linear function with
    slopes exceeding $n$ in magnitude.
    \item Conclude that $\bigcup_{n=1}^\infty A_n$ is meager, so the set of
    functions that are somewhere differentiable is meager. The ``typical''
    continuous function is nowhere differentiable.
\end{enumerate}
\end{exercise}

%% ============================================================
%% SECTION 9: SUMMARY
%% ============================================================
\section{Summary of Key Results}
\label{sec:summary}

\begin{center}
\renewcommand{\arraystretch}{1.4}
\begin{tabular}{@{}p{4.2cm}p{5.5cm}p{4.5cm}@{}}
\toprule
\textbf{Result} & \textbf{Statement (Informal)} & \textbf{RL/ML Relevance} \\
\midrule
Metric space axioms (Def.~\ref{def:metric})
  & Distance function with positivity, symmetry, triangle inequality
  & Foundation for convergence analysis \\
Completeness of $\R^n$ (Thm.~\ref{thm:Rn_complete})
  & Every Cauchy sequence in $\R^n$ converges
  & Finite-state value iteration converges \\
Completion theorem (Thm.~\ref{thm:completion})
  & Every metric space embeds densely in a complete space
  & Justifies function approximation in incomplete spaces \\
Norm equivalence in $\R^n$ (Thm.~\ref{thm:norm_equiv})
  & All norms on $\R^n$ are equivalent
  & Convergence in finite-state RL is norm-independent \\
Banach = complete normed (Def.~\ref{def:banach})
  & Banach spaces are the natural setting for fixed-point theorems
  & Value function spaces must be Banach for contraction arguments \\
Series characterization (Thm.~\ref{thm:banach_series})
  & Absolute convergence $\Rightarrow$ convergence iff Banach
  & Practical completeness test \\
$C[a,b]$ is Banach under $\|\cdot\|_\infty$ (Thm.~\ref{thm:Cab_banach})
  & Uniform limit of continuous functions is continuous
  & Value functions on compact state spaces form a Banach space \\
Bounded $\Leftrightarrow$ continuous (Thm.~\ref{thm:bounded_iff_cts})
  & A linear operator is bounded iff continuous
  & Bellman operator continuity \\
$B(X,Y)$ Banach when $Y$ Banach (Thm.~\ref{thm:BXY_banach})
  & Operator space inherits completeness
  & Space of operators on value functions is complete \\
Dual is always Banach (Cor.~\ref{cor:dual_banach})
  & $X^*$ is Banach even if $X$ is not
  & Policy gradient as element of dual space \\
Compactness equivalence (Thm.~\ref{thm:compact_equiv})
  & Compact $\Leftrightarrow$ seq.\ compact $\Leftrightarrow$ complete + totally bounded
  & Existence of optimal policies on compact spaces \\
Riesz's theorem (Thm.~\ref{thm:riesz_compact})
  & Closed unit ball compact $\Leftrightarrow$ finite-dimensional
  & Explains challenges in infinite-dim.\ function approximation \\
Arzel\`{a}--Ascoli (Thm.~\ref{thm:arzela_ascoli})
  & Equicontinuous + bounded $\Leftrightarrow$ relatively compact
  & Convergence of Lipschitz value function sequences \\
Baire category (Thm.~\ref{thm:baire})
  & Complete metric space $\ne$ countable union of nowhere dense sets
  & Underpins Uniform Boundedness Principle \\
\bottomrule
\end{tabular}
\end{center}

\bigskip

\begin{rlconnection}[Looking Ahead]
With metric spaces, norms, completeness, operators, and compactness in hand, we are
ready for the next lesson, which covers:
\begin{itemize}[nosep]
    \item The \textbf{Banach Fixed-Point Theorem} (contraction mapping theorem) ---
    the core tool for proving convergence of value iteration, policy iteration, and
    TD-learning.
    \item The \textbf{Uniform Boundedness Principle} --- used to establish uniform
    error bounds in function approximation.
    \item The \textbf{Open Mapping Theorem} and \textbf{Closed Graph Theorem} ---
    structural results for bounded operators between Banach spaces.
    \item \textbf{Hilbert spaces} --- inner product spaces where the Projection
    Theorem enables least-squares value function approximation.
\end{itemize}
\end{rlconnection}

\end{document}
